\section{Covariance Under Independent PMU Noise}
\label{app:oracle-pmu-covariance}
This retrospective replay uses previously generated physical trajectories
and nominal operating points from a split with disjoint calibration and test parents.
Per split, four event records contain 12 load and 12 line interventions.
Two normal records contribute $0.3267$ monitored hours after startup.
Each arm selects its threshold on calibration records at a common
target of one normal alarm per hour.

The replay adds independent polar-coordinate Gaussian innovations to
the original observations. The scales in the table are standard
deviations of relative magnitude and angle in radians, expressed as
percent. Magnitude scales use a fixed startup reference. The voltage
angle perturbations also enter frequency and ROCOF through the causal
21-report quadratic derivative. The resulting noise is correlated
across these channels and across successive reports.

We compare six prefix-estimated family variances, a scaled
Ledoit--Wolf covariance, a structured estimator with eight added
noise scales fitted from startup voltage-angle differences, and an
oracle that receives those eight scales from the generator.
The last two arms augment the bank with the voltage-angle noise
history. The successor-state gain includes the correlation between
process and observation noise generated by the current shock.
All arms rebuild the gains and likelihoods. The physical operators,
startup intercept, original boundary proposal scaling and publication
clock follow the same rules. The oracle specifies the added noise
component. The background noise and model discrepancy retain
their original approximations.

\begin{table}[t]\centering\small
\caption{Independent PMU noise replay on previously generated test parents.
Load and line entries count correct attribution before restoration.
Extras are unmatched alarms in event records. Normal is the alarm count
in $0.3267$ monitored hours for each arm.}
\label{tab:oracle-pmu-covariance}
\setlength{\tabcolsep}{3.4pt}
\begin{tabular}{@{}llrrrr@{}}\toprule
Scale (\%) & Covariance & Load & Line & Extras & Normal\\\midrule
0.05 & Family scale & 11/12 & 12/12 & 0 & 0 \\
0.05 & Shrinkage & 10/12 & 11/12 & 1 & 0 \\
0.05 & Structured & 11/12 & 12/12 & 0 & 0 \\
0.05 & Oracle & 11/12 & 11/12 & 1 & 0 \\
0.1 & Family scale & 11/12 & 12/12 & 0 & 0 \\
0.1 & Shrinkage & 11/12 & 12/12 & 0 & 0 \\
0.1 & Structured & 9/12 & 10/12 & 3 & 0 \\
0.1 & Oracle & 10/12 & 10/12 & 4 & 0 \\
\bottomrule
\end{tabular}

\end{table}

Even zero normal alarms in this test exposure give a one-sided 95\%
Poisson upper bound of approximately $9.17$ alarms per hour. This
replay measures sensitivity to covariance under a controlled added-noise
law. The family-scale estimator has the highest combined timely count at
both scales (tied with the structured estimator at the lower scale and
shrinkage at the higher scale). Supplying the true added-noise scales
did not improve on the family-scale arm in either replay. This comparison
does not supply the true background covariance or remove predictor error,
and the measured normal exposure cannot establish the calibration target
in test.
